\section{결론}
\label{sec:concl}

본 논문에서는 한국형 양자내성암호 서명 후보인 \texttt{HAETAE}의 응답 연산
$\zvec=\yvec+(-1)^b c\svec_1$에 대한 단일 오류주입공격을 분석하고, $+\yvec$ 덧셈을 건너뛰는
경우(T2) 단 한 번의 오류만으로 비밀키 $\svec_1$을 직접 복원할 수 있음을 STM32F405 전체 서명
환경에서 \textbf{오류 모델 하에} 실증하였다. 또한 이를
효율적으로 방어하기 위해, 연산 재계산 검증과 서명 경계 노름
재검사, 비밀키 무결성 검사를 \textbf{비교 연산에 의존하지 않는 감염형 마스킹}으로 통합한 경량 대응
기법 \textbf{IRV}를 제안하였다. 제안 기법은 거부 우회를 포함한 7개 오류 지점을 모두
차단하고(탐지 분기 스킵 모델 하) 2차 오류에 내성을 가지며, 이중 연산의 절반 미만 \emph{서명 시간}으로
동작함을 실험적으로 입증하였다. 이로써 IRV는 서명 시간과 안전성(거부 우회 차단$\cdot$2차 오류 내성)
측면에서 이중 연산 및 서명 후 검증보다 유리하다(코드$\cdot$정적 RAM은 이중 연산이 더 작다). 향후
연구로는 T1 차분(서명 2개) 공격과 동일 구조를 갖는 다른 격자 서명으로의 확장을 계획한다.
